%%%%%%%%%%%%%%%%%%%%%%%%%%%%%%%%%%%%%%%%%%%%%%%%%%%%%%%%%%%%%%%%%%%%%%%%
% 1. 常用数学符号简写
%%%%%%%%%%%%%%%%%%%%%%%%%%%%%%%%%%%%%%%%%%%%%%%%%%%%%%%%%%%%%%%%%%%%%%%%
% 希腊字母与常用符
\newcommand{\eps}{\epsilon}
\newcommand{\veps}{\varepsilon}
\newcommand{\lm}{\lambda}
\newcommand{\bs}[1]{\boldsymbol{#1}} % 粗体数学符号
\newcommand{\exs}{\exists}

% 括号与大小控制
\newcommand{\lt}{\left}
\newcommand{\rt}{\right}
\newcommand{\st}{\strut}           % 用于支撑行高
\newcommand{\dps}[1]{\displaystyle{#1}}

% 旋转符号 (用于图表或特殊表示)
\newcommand{\uin}{\mathbin{\rotatebox[origin=c]{90}{$\in$}}}
\newcommand{\usubset}{\mathbin{\rotatebox[origin=c]{90}{$\subset$}}}

% 文本缩写
\newcommand{\iid}{\textit{i.i.d.\ }} % 独立同分布，注意末尾空格

%%%%%%%%%%%%%%%%%%%%%%%%%%%%%%%%%%%%%%%%%%%%%%%%%%%%%%%%%%%%%%%%%%%%%%%%
% 2. 数学算子 (Operator) - 自动处理间距
%%%%%%%%%%%%%%%%%%%%%%%%%%%%%%%%%%%%%%%%%%%%%%%%%%%%%%%%%%%%%%%%%%%%%%%%
% 使用 DeclareMathOperator 可以在算子和变量间自动加空格
\DeclareMathOperator{\lub}{\text{lub}}
\DeclareMathOperator{\id}{\text{id}}

% 范数 Norm (依赖 mathtools)
% 用法: \norm{x} 或 \norm*{x} (自动调整大小)
\DeclarePairedDelimiter{\norm}{\lVert}{\rVert}
\newcommand{\inorm}{\norm[\infty]} % 无穷范数

%%%%%%%%%%%%%%%%%%%%%%%%%%%%%%%%%%%%%%%%%%%%%%%%%%%%%%%%%%%%%%%%%%%%%%%%
% 3. 微积分与导数
%%%%%%%%%%%%%%%%%%%%%%%%%%%%%%%%%%%%%%%%%%%%%%%%%%%%%%%%%%%%%%%%%%%%%%%%
% 偏导数 (Partial Derivatives)
\newcommand{\del}[2]{\frac{\partial #1}{\partial #2}}
% 修正：高阶导数分母通常是 \partial x^n
\newcommand{\Del}[3]{\frac{\partial^{#1} #2}{\partial #3^{#1}}}

% Display style 偏导数 (使用 \dfrac)
\newcommand{\deld}[2]{\dfrac{\partial #1}{\partial #2}}
\newcommand{\Deld}[3]{\dfrac{\partial^{#1} #2}{\partial #3^{#1}}}

% 常微分 (Ordinary Derivatives)
\newcommand{\der}[2]{\frac{\mathrm{d} #1}{\mathrm{d} #2}}
\newcommand{\ddd}[3]{\frac{\mathrm{d}^{#3} #1}{\mathrm{d} #2^{#3}}}

%%%%%%%%%%%%%%%%%%%%%%%%%%%%%%%%%%%%%%%%%%%%%%%%%%%%%%%%%%%%%%%%%%%%%%%%
% 4. 集合与拓扑专用
%%%%%%%%%%%%%%%%%%%%%%%%%%%%%%%%%%%%%%%%%%%%%%%%%%%%%%%%%%%%%%%%%%%%%%%%
\newcommand{\opensets}{\{V_{\alpha}\}_{\alpha\in I}}
\newcommand{\oset}{V_{\alpha}}
\newcommand{\opset}[1]{V_{\alpha_{#1}}}

%%%%%%%%%%%%%%%%%%%%%%%%%%%%%%%%%%%%%%%%%%%%%%%%%%%%%%%%%%%%%%%%%%%%%%%%
% 5. 格式化与解题环境
%%%%%%%%%%%%%%%%%%%%%%%%%%%%%%%%%%%%%%%%%%%%%%%%%%%%%%%%%%%%%%%%%%%%%%%%
% 段落缩进控制
\newcommand{\parinn}{\setlength{\parindent}{1cm}}
\newcommand{\parinf}{\setlength{\parindent}{0cm}}

% 解题格式
% 注意：建议改用 tcolorbox 或 proof 环境，但保留此命令以兼容旧习惯
\newcommand{\Qed}{\hfill\qed} % 简化了原本的 flushright

\newcommand{\sol}{%
    \par\setlength{\parindent}{0cm}%
    \textbf{\textit{Solution:}}\space%
    \setlength{\parindent}{1cm}%
}

\newcommand{\solve}[1]{%
    \par\setlength{\parindent}{0cm}%
    \textbf{\textit{Solution: }}\space%
    \setlength{\parindent}{1cm}#1 \Qed%
}
